在问题三主策略中，9月23日紧急购电量为170.93 kWh、实际电费为46315.30元；3月20日紧急购电量为813.93 kWh、实际电费为44999.49元。不同日期的更新收益还取决于负荷、光伏、库存与调整价格，故总体经济性以“全部更新—仅0时”完整334天控制实验评价。

购电表斜线两侧分别表示午夜原计划与最终有效购电量，二者之差为最终调整量。储能表记录实际充放电量，紧急购电表记录安全执行后仍需补足的电量，三类结果共同构成当日实际供电账本。
